%============================= app2.tex================================
\chapter{Closed-Form Minimum Distance Between Line Segments}
\label{app:distance}

Chapter~\ref{chap:distance} states the closed-form minimum-distance formulas for parallel, skew, and
intersecting segment pairs (Equations~\eqref{eq:dm_parallel}--\eqref{eq:dm_skew}) directly, without
deriving them from the underlying optimization problem. This appendix supplies that derivation. Using
the notation of Section~\ref{sec:dm_classification}, segment $L_1$ is parametrized $p(s) = a_1 +
s\mathbf{u}$, $s\in[0,1]$, $\mathbf{u} = a_2-a_1$, and segment $L_2$ is parametrized $q(t) = b_1 +
t\mathbf{v}$, $t\in[0,1]$, $\mathbf{v} = b_2-b_1$, with $\mathbf{w} = b_1-a_1$ and $\mathbf{c} =
\mathbf{u}\times\mathbf{v}$ as in the main text.

%% ============================================================
\section{The Unconstrained Problem and the Skew Case}\label{sec:appB_skew}
%% ============================================================

The squared distance between a point on each segment is
\begin{equation}
f(s,t) = \lVert p(s)-q(t) \rVert^2 = \lVert -\mathbf{w} + s\mathbf{u} - t\mathbf{v} \rVert^2,
\label{eq:appB_objective}
\end{equation}
using $p(s)-q(t) = (a_1-b_1) + s\mathbf{u} - t\mathbf{v} = -\mathbf{w} + s\mathbf{u} - t\mathbf{v}$.
Minimizing first over $s,t \in \mathbb{R}$ (ignoring the segment endpoints for the moment) by
setting the partial derivatives to zero,
\begin{equation}
\frac{\partial f}{\partial s} = 2\mathbf{u}\cdot(-\mathbf{w}+s\mathbf{u}-t\mathbf{v}) = 0, \qquad
\frac{\partial f}{\partial t} = -2\mathbf{v}\cdot(-\mathbf{w}+s\mathbf{u}-t\mathbf{v}) = 0,
\label{eq:appB_gradzero}
\end{equation}
and collecting the $s,t$ terms gives the normal equations
\begin{equation}
(\mathbf{u}\cdot\mathbf{u})s - (\mathbf{u}\cdot\mathbf{v})t = \mathbf{u}\cdot\mathbf{w}, \qquad
-(\mathbf{u}\cdot\mathbf{v})s + (\mathbf{v}\cdot\mathbf{v})t = -\mathbf{v}\cdot\mathbf{w},
\label{eq:appB_normaleqs}
\end{equation}
which is exactly the $2\times2$ linear system of Equation~\eqref{eq:dm_skew} in matrix form. Its
coefficient matrix is the Gram matrix of $\mathbf{u}$ and $\mathbf{v}$, positive semi-definite by
construction, and singular exactly when $\mathbf{u}$ and $\mathbf{v}$ are parallel, since the
determinant $(\mathbf{u}\cdot\mathbf{u})(\mathbf{v}\cdot\mathbf{v}) - (\mathbf{u}\cdot\mathbf{v})^2$
is zero, by the equality case of the Cauchy--Schwarz inequality, if and only if $\mathbf{v} =
\lambda\mathbf{u}$ for some scalar $\lambda$. For the skew case, $\lVert\mathbf{c}\rVert \geq
\varepsilon$ by definition, so $\mathbf{u}$ and $\mathbf{v}$ are not parallel, the Gram matrix is
invertible, and Equation~\eqref{eq:appB_normaleqs} has the unique solution $(s^*, t^*)$ of Equation
\eqref{eq:dm_skew}. Because $f$ is a convex quadratic (its Hessian is twice the same positive
semi-definite Gram matrix), this critical point is the global minimizer of the unconstrained problem.

If $s^*, t^* \in [0,1]$, the unconstrained minimizer already lies on both finite segments and is
therefore also the constrained minimizer, giving $d_S(L_1,L_2) = \lVert p(s^*)-q(t^*) \rVert$. If
either coordinate falls outside $[0,1]$, convexity of $f$ over the box $[0,1]^2$ implies the
constrained minimum lies on the box's boundary, that is, with $s$ or $t$ fixed at $0$ or $1$; fixing
one parameter reduces the two-variable problem to the one-variable problem of minimizing the distance
from a fixed endpoint of one segment to every point on the other, a point-to-segment distance. This
is exactly the endpoint-distance fallback used in Equation~\eqref{eq:dm_skew} and, by the identical
argument, in the intersecting case of Section~\ref{sec:appB_intersect} below.

%% ============================================================
\section{The Parallel Case}\label{sec:appB_parallel}
%% ============================================================

When $\lVert\mathbf{c}\rVert < \varepsilon$, $\mathbf{v} = \lambda\mathbf{u}$ and the normal-equation
system of Equation~\eqref{eq:appB_normaleqs} is singular: the two infinite lines have a continuum of
equally close point pairs rather than a unique one, since sliding along both lines simultaneously by
a matched amount leaves the separation unchanged. The distance between the two infinite lines is
instead the component of $\mathbf{w}$ perpendicular to the shared direction $\mathbf{u}$,
\begin{equation}
d_P^{\infty} = \lVert \operatorname{proj}_{\perp}(\mathbf{w}) \rVert, \qquad
\operatorname{proj}_{\perp}(\mathbf{w}) = \mathbf{w} - (\mathbf{w}\cdot\hat{\mathbf{u}})\hat{\mathbf{u}},
\qquad \hat{\mathbf{u}} = \frac{\mathbf{u}}{\lVert\mathbf{u}\rVert},
\label{eq:appB_projperp}
\end{equation}
obtained by subtracting from $\mathbf{w}$ its component along $\hat{\mathbf{u}}$, leaving only the
orthogonal remainder; since $\mathbf{v} \parallel \mathbf{u}$, this remainder is simultaneously
perpendicular to $\mathbf{v}$, so Equation~\eqref{eq:appB_projperp} is the distance between the two
lines regardless of which of the two directions it is measured against. This is exactly
Equation~\eqref{eq:dm_parallel}'s $\operatorname{proj}_{\perp}(\mathbf{w})$ term. Because the two
segments have finite extent, this infinite-line distance is only realized if the segments' projected
ranges overlap; the minimum over the actual finite segments is the smaller of
Equation~\eqref{eq:appB_projperp} and the four endpoint distances, matching the $\min\{\cdot\}$
construction of Equation~\eqref{eq:dm_parallel}.

%% ============================================================
\section{The Intersecting Case}\label{sec:appB_intersect}
%% ============================================================

When $\mathbf{u}$ and $\mathbf{v}$ are not parallel, $\lVert\mathbf{c}\rVert \geq \varepsilon$, the
two infinite lines intersect if and only if they are coplanar, that is, $\mathbf{w}$ lies in the
plane spanned by $\mathbf{u}$ and $\mathbf{v}$. The general test for coplanarity of three vectors is
the scalar triple product,
\begin{equation}
[\mathbf{u},\mathbf{v},\mathbf{w}] = \mathbf{u}\cdot(\mathbf{v}\times\mathbf{w})
= \mathbf{w}\cdot(\mathbf{u}\times\mathbf{v}) = \mathbf{c}\cdot\mathbf{w},
\label{eq:appB_triple}
\end{equation}
using the cyclic identity for the scalar triple product, vanishing exactly when the three vectors
are coplanar.

The classification test actually used throughout this thesis (Section~\ref{sec:dm_classification}
and Algorithm~\ref{alg:dm_collision}) is instead $\lVert\mathbf{c}\times\mathbf{w}\rVert <
\varepsilon$, a vector rather than a scalar condition. Expanding it with the vector triple-product
identity $(\mathbf{A}\times\mathbf{B})\times\mathbf{C} = \mathbf{B}(\mathbf{A}\cdot\mathbf{C}) -
\mathbf{A}(\mathbf{B}\cdot\mathbf{C})$,
\begin{equation}
\mathbf{c}\times\mathbf{w} = (\mathbf{u}\times\mathbf{v})\times\mathbf{w}
= \mathbf{v}(\mathbf{u}\cdot\mathbf{w}) - \mathbf{u}(\mathbf{v}\cdot\mathbf{w}),
\label{eq:appB_cwtriple}
\end{equation}
which, since $\mathbf{u}$ and $\mathbf{v}$ are linearly independent here, vanishes if and only if
both scalar coefficients vanish simultaneously: $\mathbf{u}\cdot\mathbf{w} = 0$ and
$\mathbf{v}\cdot\mathbf{w} = 0$. This is the condition the implemented test actually detects, the
connecting vector $\mathbf{w}$ between the two segments' start points is perpendicular to both
direction vectors at once, a stricter condition than the general coplanarity test of Equation
\eqref{eq:appB_triple}, and the one this thesis's classification and the accompanying
implementation~\citep{sk2025distance} use throughout. Given coplanarity under either test, the
intersection parameters solve $\mathbf{u}s - \mathbf{v}t = \mathbf{w}$ as in Equation
\eqref{eq:dm_intersect}; if $s,t \in [0,1]$ the segments themselves intersect and
$d_I(L_1,L_2)=0$, otherwise the nearest approach is again an endpoint distance, by the same
boundary argument as Section~\ref{sec:appB_skew}.

%========================================================================

%%% Local  Variables:
%%% mode:  latex
%%% TeX-master: "Thesis.tex"
%%% End:
